% !TEX root =  master.tex
\chapter{Nachhaltigkeit in der Materialentwicklung}

In den vergangenen Jahren haben zunehmende Umwelt‑ und Klimaprobleme sowie Herausforderungen wie Armut, wachsende soziale Ungleichheiten und gesellschaftliche Spannungen das Thema der nachhaltigen Entwicklung verstärkt in den Fokus gerückt. Nationale und internationale Institutionen, politische Entscheidungsträger sowie praxis‑ und forschungsnahe Akteure haben der sozialen und ökologischen Nachhaltigkeit weltweit erhöhte Aufmerksamkeit gewidmet. \cite{battilana2010building}


\subsection{Nachhaltigkeit}
Der Begriff „nachhaltige Entwicklung“ beschreibt die Idee, die Bedürfnisse der Gegenwart zu befriedigen, ohne die Chancen künftiger Generationen zu beeinträchtigen. Allerdings ist der Begriff „nachhaltig“ selbst schwer präzise zu definieren: Er umfasst wirtschaftliches Wohlergehen, soziale Gerechtigkeit und die Reduzierung von Umweltschäden. Diese Ziele sind nicht immer vereinbar. Dies ist besonders problematisch, wenn Umweltschäden pro Einheit wirtschaftlicher oder physischer Leistung bewertet werden: Steigt das Produktionsvolumen schneller als die Schäden pro Produktionseinheit sinken, nehmen die gesamten Umweltschäden zu. Es sind die absoluten, nicht die relativen, physischen Folgen, die eine Bedrohung für zukünftige Generationen darstellen.


Trotz eines breiten Konsenses über die Bedeutung nachhaltiger Entwicklung wird deren konkrete inhaltliche Ausgestaltung und Zielsetzung jedoch nur selten explizit diskutiert und systematisch analysiert. Infolgedessen besteht die Gefahr, dass die praktische Umsetzung von Nachhaltigkeit durch die Unschärfe und Allgegenwärtigkeit des Begriffs eingeschränkt wird. \cite{gray2010accounting}

Aufgrund der weit verbreiteten Verwendung des Nachhaltigkeitsbegriffs haben sich im Laufe der Zeit unterschiedliche Interpretationsansätze herausgebildet, die Nachhaltigkeit beispielsweise mit sozialer Verantwortung, Umweltmanagement oder unternehmerischer Nachhaltigkeit verknüpfen und häufig getrennt voneinander betrachten. Diese unterschiedlichen Perspektiven haben zugleich die Aufmerksamkeit auf die Rolle von Unternehmen gelenkt, deren Verantwortung und Beitrag zur nachhaltigen Entwicklung zunehmend hinterfragt werden. \cite{joseph2012ambiguous}

\subsection{Nachhaltigkeit der Materialien}
Angesichts des zunehmenden Drucks durch nationale und internationale Regulierungen sowie durch gesellschaftliche Erwartungen werden Unternehmen zunehmend dazu verpflichtet, Prinzipien sozialer und ökologischer Verantwortung systematisch in ihre Strategien, Organisationsstrukturen und Managementsysteme zu integrieren. \cite{gond2012configuring}

Insbesondere die jüngsten europäischen Regulierungen, wie die Ökodesign‑Verordnung für nachhaltige Produkte (ESPR), erfordern eine deutlich beschleunigte Transformation hin zu nachhaltigen Produkten und Prozessen. Diese Anforderungen lassen sich jedoch ohne zusätzliche technologische Hilfsmittel, insbesondere ohne den Einsatz datengetriebener Methoden und Künstlicher Intelligenz, nur schwer erfüllen. \cite{faraca2024ecodesign}


Gleichzeitig sind Entwicklungszyklen für neue Materialien traditionell sehr zeit‑ und kostenintensiv und können mehrere Jahre in Anspruch nehmen, insbesondere wenn sie auf klassischen experimentellen und iterativen Entwicklungsansätzen basieren. Diese langen Entwicklungszeiten stehen zunehmend im Spannungsfeld zu regulatorischen Anforderungen und verkürzten Innovationszyklen, was den Bedarf an daten‑ und KI‑gestützten Methoden zur Beschleunigung der Materialentwicklung weiter erhöht. \cite{jakobsen2025current}


\cite{akafuah2016evolution}
\cite{tisza2013advanced}


\subsubsection{Lebenszyklusbetrachtung (LCA)}


